% !TeX root = ../main.tex
%===========================================================================%
%->> Appendix B (口試委員會提問與答辯紀錄)
%===========================================================================%
\chapter{口試委員會提問與答辯紀錄 (Oral Defense Q\&A Record)}
\label{appendix:B}

本附錄整理自本學術論文於學位口試（Oral Defense）期間，口試委員會全體委員之關鍵學術提問、審查意見，以及本研究之詳細答辯與物理機制澄清。

\section{系統建模與技術選型答辯 (System Modeling \& Tech Selection)}

\subsection*{Q1.1 [潘仁義 教授]：三種的天線域關閉天線的方式（Type 1 / Type 2 / Type 3），本研究是使用哪一種？}
\textbf{答辯與澄清}：\\
本研究採用的是 \textbf{SDA Type 1 Adaptation}（天線埠直接靜默機制）來對啟用的天線數量進行動態調整。

\subsection*{Q1.2 [李昌明 教授]：您說論文是使用 Type 1 的 Adaptation，有什麼原因嗎？看論文中好像沒有明確說明其物理原因。}
\textbf{答辯與澄清}：\\
本研究選擇採用 Type 1 天線靜默機制，主要基於以下三點學術與通訊工程考量：
\begin{enumerate}
    \item \textbf{業界研究一致性}：多數關於空間域自適應（Spatial Domain Adaptation, SDA）的文獻與 3GPP 討論皆以 Type 1 作為天線縮減的基準方案。
    \item \textbf{量測開銷（CSI Overhead）考量}：相較於 Type 2 等其他縮減方式，Type 1 能夠將降維後的通道量測直接「映射」在同一個全維度通道矩陣中。若採用 Type 2，使用者設備（UE）在通道量測上會面臨巨大的信令開銷，甚至需要基地台另外配置獨立的下行資源供其量測。
    \item \textbf{預編碼器（Precoder）兼容性}：依據 3GPP TR 38.864 標準中的 Set 1 規範，基地台採用的是 \textbf{Non-codebook based 預編碼器}。在此物理假設下，搭配 Type 1 埠數縮減（16T/32T/64T）在模擬器底層的實作最具工程可行性。
\end{enumerate}

\subsection*{Q1.3 [李昌明 教授]：天線域適應實作的部分，有沒有考慮到參考訊號（Reference Signal）可能要重新設計？}
\textbf{答辯與澄清}：\\
在本研究目前的實作範疇中，\textbf{並未重新設計參考訊號}。我們完全遵循 3GPP Release 18 標準所規規定之 CSI-RS（通道狀態資訊參考訊號）的實體層資源對齊與配置規範。

\subsection*{Q1.4 [李昌明 教授]：基地台天線關掉之後，你要怎麼讓使用者端（UE）知道當下的天線配置？這是否會花費一些控制開銷（Overhead）？如果切換得太頻繁，這個 Overhead 會不會影響到最後實際上的節能效果？}
\textbf{答辯與澄清}：\\
是的，口試委員的考量完全切中實務通訊系統的痛點。
\begin{enumerate}
    \item \textbf{信令告知機制}：目前 3GPP 標準尚未明定基地台在動態調整天線時，應如何即時告知 UE。但在現有前瞻研究中，常見的作法是在每一個下行時隙（Time Slot / TTI）的最前方加入一個專屬的下行控制資訊（Downlink Control Information, DCI），用以攜帶當下的天線決策狀態。
    \item \textbf{開銷與節能折衷}：天線組態的頻繁切換確實會引入控制面開銷（Control-plane Overhead），若切換過於盲目與頻繁，該開銷會大幅折損實質節能效益。這正是本研究\textbf{不採用傳統無緩衝 Heuristic 方案}，而是引入李雅普諾夫最佳化理論並結合\textbf{「時序延遲緩衝（Transition Delay Buffer）」}與\textbf{「非線性調控因子 $\gamma = 2.0$」}的核心原因——旨在自演算法底層徹底抑制乒乓效應，在保障硬體穩定的同時，將信令開銷控制在合理區間內。
\end{enumerate}

\section{李雅普諾夫最佳化演算法答辯 (Lyapunov Optimization Algorithm)}

\subsection*{Q2.1 [李光啟 教授]：公式 (3-14) 漂移加懲罰（Drift-Plus-Penalty, DPP）框架中，控制權重 $V$ 設定較大時是省電優先，而 $V$ 設定較小時則是效能（時延）優先，對嗎？}
\textbf{答辯與澄清}：\\
\textbf{是的，您的理解完全正確。}\\
在李雅普諾夫 DPP 框架中，我們的最佳化目標是最小化以下目標函數：$\Delta(t) + V \cdot E[P_{\text{total}}(t) \mid Q(t)]$（此處已採用標準期望值算子 $E$ 進行最廣義之學術相容表記）。
\begin{itemize}
    \item \textbf{當 $V$ 值設定極小（$V \to 0$）時}：功耗懲罰項被弱化，最佳化引擎會全力最小化李雅普諾夫漂移 $\Delta(t)$（即壓制佇列長度），系統會導向\textbf{性能（低時延）優先}。
    \item \textbf{當 $V$ 值設定較大時}：基地台功耗懲罰被放大，系統會為了最小化功耗而主動容忍封包在緩衝區內適度積壓，從而導向\textbf{節能（省電）優先}。
\end{itemize}

\subsection*{Q2.2 [李昌明 教授]：公式 (3-23) 中，平均抵達速率 $A_{\text{avg}}(t)$ 計算所採用的滑動視窗（Sliding Window）大小是 Spec 設定的嗎？這會影響到演算法的敏感度嗎？}
\textbf{答辯與澄清}：\\
該滑動視窗的大小\textbf{並非標準（Spec）定義，而是由本研究自行設計與調優的}。\\
我們在 MATLAB 演算法沙盒驗證階段，針對不同的滑動視窗大小進行了敏感度評估。實驗證實，視窗大小確實會直接影響演算法對即時隨機流量的敏感度：視窗過大會導致天線切換反應過於遲緩，失去即時反饋優勢；視窗過小則會過度放大瞬時噪訊，引發乒乓切換。本研究最終設定之視窗長度，是在兩者之間取得最佳折衷的實體物理常數。

\subsection*{Q2.3 [李光啟 教授]：為什麼會選擇使用 Lyapunov 理論做天線決策？是 3GPP 規定的嗎？還是你們自創的？除了這以外還有別的演算法嗎？}
\textbf{答辯與澄清}：\\
\begin{enumerate}
    \item \textbf{3GPP 規範與演算法邊界}：3GPP 國際標準僅定義了實體層天線域適應的硬體調整方式（如 16T/32T/64T），但\textbf{從未明定如何進行決策最佳化}，這完全留給廠商或研究者進行演算法實作。
    \item \textbf{為什麼選擇 Lyapunov}：李雅普諾夫漂移加懲罰理論是動態隨機網路控制領域的權威數學工具。相較於其他方案（如強化學習），它最大的科學優勢在於\textbf{「無需任何先驗流量知識（Requires no future traffic knowledge）」}，僅憑當前即時觀測值（Instantaneous queue \& channel states）即可在每個時隙做出最佳線上決策，且能在數學上嚴格證明佇列時延的漸進安全邊界。
    \item \textbf{系統部署可行性}：在系統級模擬（SLS）的封閉且高時效性（TTI = 0.5 ms / 1 ms）環境下，Lyapunov 演算法屬於離線/線上皆極低複雜度的解析決策方程式，避免了在模擬器中嵌入龐大機器學習或決策引擎帶來的額外運算開銷與收斂不穩定風險，極具工業落地價值。
\end{enumerate}

\subsection*{Q2.4 [許正欣 教授]：為何最佳化目標方程式 (3-26) 中，會出現「資料傳輸速率（Rate）」乘上「佇列長度（Queue size）」的項？}
\textbf{答辯與澄清}：\\
這一項是李雅普諾夫漂移加懲罰（Drift-Plus-Penalty, DPP）框架在數學上進行「條件期望值上界推導與化簡」後的必然數學產物。\\
其物理本質在於將當前 MAC 層的佇列積壓 $Q_n(t)$，作為實體層可達傳輸速率 $R_n(a(t))$ 的\textbf{隨機動態權重加權值}。這使得決策引擎在求解最小化問題時，能自適應進行物理權衡：
\begin{itemize}
    \item 當佇列長度 $Q(t)$ 極大（網路壅塞）時，該加權項會遠大於功耗項，強迫基地台啟用 64T 窄波束全力排空佇列。
    \item 當佇列長度 $Q(t)$ 趨近於零（網路閒置）時，功耗項主導決策，基地台得以退回 16T 低功耗休眠狀態。
\end{itemize}

\subsection*{Q2.5 [李光啟 教授]：關於流量感知指標 $\rho(t)$，為什麼您在公式中將其設計為非線性的 $\rho(t)^{\gamma}$ 次方，而不是常見的 S 曲線（Sigmoid）？目前 $\gamma$ 設定為 2.0，在負載約 0.5 時，感覺上演算法會過度偏向省電？}
\textbf{答辯與澄清}：\\
\textbf{這確實是本演算法在實務設計上的一個重要物理折衷（Trade-off）。}\\
\begin{enumerate}
    \item \textbf{為什麼設定 $\gamma = 2.0$}：根據我們進行的參數敏感度分析，當 $\gamma = 2.0$ 時，系統能在長期平均節能率（ESG）與用戶速率之間取得對數折衷的最佳平衡點。
    \item \textbf{省電偏好與缺陷澄清}：當 $\gamma = 2.0$ 且即時負載率為 $\rho(t) = 0.5$ 時，非線性映射的權重確實會顯著偏向最大化省電。這正是本機制能榨出高達 30.3\% 的巨額節能率（遠優於傳統方案之 21.2\%）的數學祕密。其物理代價是使用者的感知吞吐量（UPT）會產生約 1.6\% 的極輕微速率下滑。
    \item \textbf{S 曲線與未來工作}：您提出的 S 曲線（Sigmoid）或雙曲正切（tanh）函數是一個非常精彩的改進思路。S 曲線能提供中等負載下更為平滑且對稱的轉換，避免過早鎖死在省電組態。我們已將此項「將 $\gamma$ 升級為干擾與負載自適應之非線性映射器」明確列入論文最後的 Future Work（未來展望第一點）中，作為後續研究的突破口。
\end{enumerate}

\section{模擬結果與物理現象分析答辯 (Simulation \& Physical Phenomena)}

\subsection*{Q3.1 [李光啟 教授]：在第四章結果中，為什麼低負載（Low Load）與輕負載（Light Load）場景下的節能率（ESG）剛好都是 30.3\%？這是否合理？}
\textbf{答辯與澄清}：\\
\textbf{這在統計上純屬精準四捨五入後的巧合。}\\
我們在 system-level 閉環模擬中，基地台在兩個場景下因為流量隨機分佈的不同，其天線進入 16T 靜默狀態的時間比例與轉渡開銷其實存在著微觀上的物理差異。但在我們將長期平均總功耗代入公式 (3-6) 計算並取小數點後第一位時，兩者的量化 ESG 數值剛好四捨五入收斂至 30.3\%，其實際未修約的原始浮點數數據仍有微幅偏差。這再次證明了我們演算法在中低負載場景下，皆能穩健且飽和地榨取高達 30\% 左右的空間域能效红利。

\subsection*{Q3.2 [李光啟 教授]：當天線頻繁收放（Ping-pong Effect）發生時，為什麼會導致佇列時延發散（Queue Delay Divergence）？這背後的通訊物理是什麼？}
\textbf{答辯與澄清}：\\
這是本論文在控制理論與排隊網路交叉領域的核心研究亮點，其物理因果關係主要由以下三個維度驅動：
\begin{enumerate}
    \item \textbf{服務速率方差之二次方放大}：依據隨機排隊論著名的 Pollaczek-Khinchine (P-K) 公式，使用者的平均排隊佇列長度除了與平均服務率相關外，更與服務速率的\textbf{統計方差（Variance $\sigma^2_\mu$）}呈正比。天線的頻繁、盲目切換會使實體層速率產生極大的瞬時震盪，這使得方差 $\sigma^2_\mu$ 呈二次方級數飆升，直接導致佇列長度與封包延遲急遽發散。
    \item \textbf{非負邊界約束下之非對稱累積}：佇列的演進具有強烈之非線性。當突發流量到來而天線因 Heuristic 反應遲鈍鎖死在 16T 時，佇列會非線性飆高；而當天線後知後覺地終於開到 64T 時，若此時佇列已被排空，超額提供的傳輸速率將會被 $\max$ 算子完全歸零（即無法補償前期的時延損失），進而產生了極為嚴重的長尾時延。
    \item 本研究透過流量感知動態 weight，徹底消除了切換抖動，成功將佇列時延安全壓制並收斂在安全邊界內。
\end{enumerate}

\subsection*{Q3.3 [李光啟 教授]：您使用了一個機制叫做 Enhance CSI，你怎麼樣去量測那些處於 OFF 狀態（已靜默關閉）的天線通道？}
\textbf{答辯與澄清}：\\
這正是本研究採用 Type 1 空間域適應 搭配 Enhanced CSI 探測週期 的物理核心精髓。
\begin{enumerate}
    \item \textbf{全維度探測時機}：處於靜默狀態（OFF）的天線在日常數據傳輸（PDSCH）時確實不發送訊號。但是，在 Enhanced CSI 觸發時隙（其探測週期設定為每 10 個時隙觸發一次），基地台會\textbf{短暫且強制地重啟所有天線埠}，並發送一次完整的、覆蓋 64 埠的全維度 CSI-RS 探測訊號進行閉環信道量測。
    \item \textbf{零遮罩推估（Zero Masking）}：當 UE 將此全維度通道矩陣反饋回基地台後，基地台決策大腦再利用我們設計的 Zero Masking 方陣進行即時降維。這使得我們\textbf{僅需發送一次探測訊號，就能完美推估出 16T、32T、64T 三種候選方案在當前無線信道下的等效 SINR 與預期傳輸速率}，成功解決了靜默天線無法量測的物理障礙。
\end{enumerate}

\section{論文格式與模擬平台細節答辯 (Thesis Format \& Simulation Details)}

\subsection*{Q4.1 [許正欣 教授]：論文中所有圖（Figure）下方的說明文字，不需要加句號。}
\textbf{答辯與澄清}：\\
\textbf{感謝口試委員的嚴謹指正。}\\
我們已對全論文的所有圖目錄與圖下方之 Caption 文字進行了拉網式排查，\textbf{已將所有說明文字末尾的句號（Period）全部刪除}，完全符合學術圖表標準排版規範。

\subsection*{Q4.2 [許正欣 教授]：通道矩陣的虛數單位 $j$，在論文多處存在大小寫標記不一致的錯誤。}
\textbf{答辯與澄清}：\\
\textbf{感謝教授的仔細審查。}\\
我們已將論文中所有涉及複數矩陣與通道公式的虛數單位，\textbf{統一修正為斜體斜寫的小寫 $j$}，消除了所有大小寫混淆的筆誤，維持了整本論文公式的高度學術一致性。

\subsection*{Q4.3 [李昌明 教授]：我記得 3GPP TR 38.864 規範中，基地台有五種不同的功耗狀態（Power State）？為什麼你們論文的表 3-2 功耗模型中只有列出兩三個？}
\textbf{答辯與澄清}：\\
\begin{enumerate}
    \item \textbf{五大狀態事實}：3GPP TR 38.864 技術報告中確實定義了五種功耗狀態：包括 $P_1$（全開）、$P_2$（部分休眠）、$P_3$（微睡眠 Micro-sleep）、$P_4$（Active DL 數據傳輸）以及 $P_5$（Active UL 接收）。
    \item \textbf{空間域自適應的實務專注}：在探討空間域自適應（SDA）天線開關的研究中，業界與標準自適應評估通常只專注在基地台\textbf{處於下行 Active 數據傳輸狀態（Active DL, 即 $P_4$）}與\textbf{待機微睡眠狀態（Micro-sleep, 即 $P_3$）} 之間的切換。因此，為避免過度偏離主題，本研究將功耗模型的建立專注於 $P_3$（基礎靜態能耗）與 $P_4$（動態射頻調適能耗）的線性物理耦合，這完全符合 TR 38.864 空間域適應評估的學術標準規範。
\end{enumerate}